%%%BLOCK id=005-02 ch=14 sec=电磁波谱 kind=summary alt=17 cont=none y=0.235-0.30
\noindent 无线电波\quad 微波\quad 红外线\,可见光\textsubscript{\,红\ 紫}\quad 紫外线\quad X射线\textsuperscript{(伦琴)}\quad $\gamma$射线\quad $\longrightarrow$\quad $f$大\ $\lambda$小 % 原稿：自左向右的长箭头画在整行文字上方，箭头末端写“f大”，其下写“λ小”；“(伦琴)”在“X射线”上方，“红 紫”在“可见光”下方
%%%END
%%%BLOCK id=105-02 ch=14 sec=麦克斯韦电磁场理论 kind=knowledge alt=none cont=none y=0.24-0.34
②\quad \textbf{麦克斯韦电磁场理论}

③\quad 变化电/磁场在周围产生变化的磁/电场 % 圈号②③位于左页边，对应关系按垂直位置推定
%%%END
%%%BLOCK id=105-03 ch=14 sec=电磁波的发射与接收 kind=knowledge alt=none cont=none y=0.35-0.53
④\quad \textbf{电磁波的发射与接收，}\quad 电视、雷达、移动电话
\begin{itemize}
\item 发射电磁波要开放的高频振荡电路.\\ 并对电磁波根据信号的强弱进行\underline{调幅}、\underline{调频} % 原稿此两行用左大括号括在一起
\item 接收电磁波需要能产生电谐振的调谐电路，再把信号从高频电流中解调出来\\ 调幅波的解调也叫\underline{检波} % 原稿此两行用左大括号括在一起
\end{itemize}
%%%END
%%%BLOCK id=106-04 ch=14 sec=电磁波与机械波对比 kind=summary alt=08 cont=none y=0.60-0.80
\begin{center}
\begin{tabular}{@{}lp{0.34\linewidth}p{0.4\linewidth}@{}}
\toprule
 & 电磁波 & 机械波 \\
\midrule
产生 & 由周期性变化的电场、磁场产生. & 由质点（波源）的振动产生 \\
波的特点. & 横波 & 纵波或横波. \\
波速 & 真空中等于光速$c=3\times10^{8}\unit{m/s}$ & 在空气中不大（声波在空气中$340\unit{m/s}$） \\
是否需要介质 & 不需要介质，可在真空中传播 & 必须有介质，真空中不能传播. \\
能量传播 & 电磁能 & 机械能. \\
\bottomrule
\end{tabular}
\end{center}
%%%END
%%%BLOCK id=106-05 ch=14 sec=麦克斯韦电磁场理论 kind=knowledge alt=none cont=none y=0.80-0.89
% 原稿此三行用左大括号括在一起
\begin{itemize}
\item 均匀变化的电场（磁场）产生 均匀变化的磁场（电场） % CHECK: 原文如此(物理上应为稳定的磁场(电场))
\item 非均匀变化的电场（磁场）产生 非均匀变化的磁场（电场）.
\item 周期性振荡的电场（磁场）产生 周期性振荡的磁场（电场）.
\end{itemize}
%%%END
%%%BLOCK id=106-06 ch=14 sec=LC振荡 kind=note alt=none cont=none y=0.89-1.00
$I\propto$磁场能\qquad $U\propto$电场能.\qquad $T$=2次充放电.

\struck{$U_c$的导数为$U_L$}

\struck{互为导数} % 原稿上述两行均被横线划去
%%%END
%%%BLOCK id=113-01 ch=14 sec=电磁波谱 kind=summary alt=none cont=none y=0.04-0.14
无线电波\quad 微波\quad 红外线\quad 可见光\quad 紫外线\quad $X$射线\quad $\gamma$射线 $\longrightarrow$\quad $f\uparrow$\quad $\lambda\downarrow$

\fbox{无微红可紫 $X$ $\gamma$}
%%%END
